\section{Behavior Evals: de «algo no está bien en el modelo» a un sistema de rechazo depurable}

El comportamiento de un modelo no es como un problema de matemáticas que tiene una sola respuesta. Los rechazos, el carácter directo, el equilibrio en temas políticos, el reconocimiento de alucinaciones, el estilo lingüístico y las reacciones emocionales conforman en conjunto el producto que percibe el usuario. La ponente pregunta primero «cómo hacer una eval confiable» y después usa la mejora del over-refusal de Claude 2.1 a Claude 3 para mostrar que el proceso correcto no es ajustar por intuición un reward total, sino construir una taxonomy, combinar fuentes de datos, rastrear el problema hasta conjuntos de entrenamiento concretos y hacer regresión con ejemplos before/after y benchmarks.

\lecturefigure{slide-34-evals-question.jpg}{Pregunta central: cómo hacer una eval confiable y cómo medir el comportamiento que realmente se necesita}{00:20:25--00:20:43}

\begin{importantbox}{Definición: una eval es un protocolo de decisión, no un archivo de preguntas}
Una eval completa contiene al menos la distribución de tareas $D$, la versión del modelo/sistema $M$, la regla de calificación $S$, las condiciones de ejecución $C$ y el umbral de aceptación $\tau$:
\[
\operatorname{Eval}(M;D,S,C)\ge \tau.
\]
Si cambian el prompt, las herramientas, el system message, el sampling, la versión del judge o la rubric humana, los resultados ya no pueden compararse directamente. Lo que vuelve confiable una eval es que sea reproducible y que pueda sustentar decisiones de producto, no la cantidad de preguntas.
\end{importantbox}

\subsection{Escribir el model behavior como un vector}

El primer grupo de traits incluye reducir el benign refusal, responder con una postura más definida y disminuir las milquetoast response, al tiempo que se conservan caveats donde hay incertidumbre. El segundo grupo de traits incluye el nuance en temas de alto riesgo, fomentar la media literacy, dar explicaciones moderadamente más largas y saber lo que uno no sabe. No forman una misma dimensión de «helpfulness»; si solo se optimiza una puntuación única, el modelo puede encubrir errores factuales con textos más largos y más complacientes.

\lecturefigure{slide-35-general-behaviors-refusal-opinions.jpg}{Vector de comportamiento I: benign refusal, carácter directo, caveat y respuestas que no siguen una plantilla}{00:20:43--00:23:23}

\lecturefigure{slide-36-general-behaviors-nuance-knowledge.jpg}{Vector de comportamiento II: nuance, media literacy, longitud y calibrated self-knowledge}{00:23:23--00:24:17}

El comportamiento de una respuesta puede representarse como
\[
\mathbf{b}(x,y)=
\begin{bmatrix}
h & s & d & n & c & v
\end{bmatrix}^{\top},
\]
donde $h$ es helpfulness, $s$ es safety, $d$ es directness, $n$ es nuance, $c$ es calibration y $v$ es voice consistency. La aceptación del producto no consiste en buscar un máximo global, sino en definir una región permitida $\mathcal{B}_{\text{task}}$ para cada dominio de riesgo. Por ejemplo, las preguntas médicas exigen caveats más fuertes, mientras que la escritura creativa no debería rechazarse automáticamente porque aparezcan palabras como «vigilancia» o «heist».

\begin{knowledgebox}{Términos clave: dimensiones comunes de la evaluación de comportamiento}
\begin{tabularx}{\textwidth}{L{0.18\textwidth} X X}
\toprule
Dimensión & Problema que resuelve & Qué no puede sustituir \\
\midrule
Helpfulness & Si hace avanzar el objetivo del usuario & No sustituye la corrección factual \\
Harmlessness & Si evita daños reales & No puede juzgar la intención por palabras clave superficiales \\
Directness & Si responde primero la pregunta central & No puede omitir las salvedades necesarias \\
Nuance & Si conserva la evidencia en conflicto y las condiciones & No puede convertirse en complacer a ambas partes sin llegar a una conclusión \\
Calibration & Si expresa un grado de confianza razonable & No sustituye la verificación externa \\
Self-knowledge & Si conoce los límites de su capacidad o de su información & No puede simularse con descargos de responsabilidad fijos \\
\bottomrule
\end{tabularx}
\end{knowledgebox}

\teachervoice{00:20:25--00:23:52, la clase no resumió «Claude 3 es mejor» con una sola puntuación, sino que enumeró uno por uno refusals, directness, nuance, balance y self-knowledge. Nota de clase: primero hay que descomponer el comportamiento del producto; solo así el equipo sabe qué aspecto mejoró y cuál se sacrificó.}

\subsection{Over-refusal: el clasificador de seguridad toma las palabras superficiales por intención}

El problema típico de Claude 2.1 es que el prompt contiene en apariencia palabras peligrosas, aunque la tarea real es inofensiva. En la figura, «kill a python process» es una pregunta de gestión de procesos; si el modelo solo asocia el riesgo por keyword, la rechazará. Más difíciles de detectar son los casos de escritura creativa, de documentos largos adjuntos y de function calling: el sistema puede producir un misdirected refusal porque algún pasaje del contexto parece dañino o porque olvida el system prompt.

\lecturefigure{slide-37-claude21-overrefusal.jpg}{Over-refusal de Claude 2.1: palabras peligrosas superficiales hacen que se rechacen solicitudes inofensivas}{00:23:52--00:25:55}

\lecturefigure{slide-38-nuanced-refusals.jpg}{Rechazos más matizados: explicar la responsabilidad, reconocer los límites y ofrecer alternativas viables}{00:25:55--00:26:27}

Sea $D_b$ el conjunto de benign prompts y $r(x)\in\{0,1\}$ el indicador de rechazo del modelo; entonces la false refusal rate es
\[
\operatorname{FRR}=\frac{1}{|D_b|}\sum_{x\in D_b}r(x).
\]
Que la FRR baje no significa automáticamente que la seguridad haya mejorado, porque el modelo también puede empezar a responder solicitudes realmente dañinas. Por eso debe reportarse al menos junto con la harmful compliance rate y estratificarse por dominio; el promedio total oculta fallas locales en function calling, creative writing o documentos largos.

\begin{warningbox}{Rechazar menos no siempre es mejor}
El objetivo de optimización debería ser «rechazar menos las solicitudes inofensivas, limitar de forma robusta los riesgos reales y ofrecer ayuda segura y explicable ante las solicitudes limítrofes». Reducir solo la FRR induce unsafe compliance; reducir solo la harmful compliance induce rechazos basados en plantillas. Ambos casos requieren hard negatives: pares de ejemplos parecidos en la superficie pero con intenciones distintas.
\end{warningbox}

\subsection{Taxonomy: primero nombrar la falla y después decidir qué reparar}

La taxonomy que enumera la ponente incluye benign over-refusal, creative-writing refusal, function-calling refusal, long-document attachment y misdirected refusal. La clasificación no sirve para redactar informes, sino para mapear «el modelo rechazó» a la capa del sistema que podría causarlo: datos de seguridad, datos de self-knowledge, tool schema, enrutamiento de long-context u olvido del system prompt.

\lecturefigure{slide-39-refusal-taxonomy.jpg}{Taxonomy de rechazos: benign, creative writing y function calling}{00:26:27--00:27:08}

\lecturefigure{slide-40-refusal-taxonomy-examples.jpg}{Ejemplos visuales de la taxonomy de rechazos: adjuntos y misdirected refusal}{00:27:08--00:28:12}

\begin{knowledgebox}{Tabla de root cause: causas distintas bajo una misma apariencia de rechazo}
\begin{tabularx}{\textwidth}{L{0.22\textwidth} X X}
\toprule
Síntoma & Posible causa raíz & Qué revisar primero \\
\midrule
Se rechaza ``kill python process'' & Datos de seguridad basados en palabras clave con demasiado peso & benign hard negatives \\
Se rechaza una historia creativa & Los ejemplos de seguridad contaminan el dominio creativo & Proporción entre dominios y pair data \\
``can't see tool'' & Datos de self-knowledge erróneos & Estado de disponibilidad de las herramientas \\
Un adjunto largo provoca el rechazo & Contenido local del documento induce un juicio global erróneo & Atribución por párrafo e intención de la consulta \\
Olvida el system prompt & Defecto de contexto o de enrutamiento & Ensamblado del prompt y trazas de regresión \\
\bottomrule
\end{tabularx}
Un mismo nombre, «refusal», puede requerir reparaciones completamente distintas; agregar primero datos genéricos de SFT suele provocar contaminación cruzada.
\end{knowledgebox}

\begin{lstlisting}[caption={Flujo mínimo de triaje para fallas de rechazo}]
case = replay(prompt, attachments, tools, system_prompt, model_version)
if case.intent_is_benign and case.refused:
    bucket = classify_refusal(case)
    sources = trace_training_and_policy_influences(bucket)
    candidate_fix = design_targeted_data_or_system_change(sources)
    run_regression(candidate_fix, benign_suite, harmful_suite, domain_suite)
else:
    preserve_or_strengthen_safety_boundary(case)
\end{lstlisting}

\teachervoice{00:23:52--00:28:35, Karina enfatiza repetidamente que el over-refusal no proviene de una sola fuente de datos; el equipo construyó una taxonomy que distingue benign, creative writing, function calling, long document y misdirected refusal. Experiencia práctica: ante un mal ejemplo, primero hay que preguntar a qué mecanismo pertenece, en lugar de modificar de inmediato el reward global.}

\subsection{Combinación de fuentes para una eval confiable}

Las slides dividen los eval prompts en fallas reales recopiladas en el product flywheel, ejemplos de usuarios curados manualmente, synthetic prompts que cubren la frontera harmless/helpful y benchmarks externos como XSTest. Cada tipo atiende un punto ciego distinto: el tráfico real tiene validez ecológica pero un muestreo sesgado; los synthetic data son controlables pero heredan el sesgo del generator; los benchmarks públicos son reproducibles pero propensos al sobreajuste.

\lecturefigure{slide-41-trusted-evals-sources.jpg}{Fuentes de una eval confiable: product flywheel, synthetic prompts y XSTest}{00:28:12--00:30:21}

Si la distribución de la $k$-ésima fuente es $D_k$, la eval combinada puede escribirse como una mixture
\[
D_{\text{eval}}=\sum_{k=1}^{K}w_kD_k,
\qquad \sum_k w_k=1.
\]
$w_k$ no es una constante natural, sino una decisión de riesgo del producto. Ponderar según la proporción del tráfico subestima los riesgos altos de la cola larga; una ponderación completamente uniforme, en cambio, puede alejarse del uso real. Lo más prudente es reportar a la vez overall, slice y worst-case, y congelar los conjuntos de regresión críticos.

\begin{knowledgebox}{Primera mención: RLAIF y synthetic preference data}
\term{RLAIF (Reinforcement Learning from AI Feedback)} usa un AI judge o preferences generadas por IA para sustituir o complementar la retroalimentación humana. Los synthetic preference data permiten cubrir con rapidez estilos, límites y casos de cola larga específicos, pero su valor depende de la diversity, de la judge calibration y de la auditoría por muestreo. No se trata de «la IA demostrando por sí misma que tiene razón», sino de un proceso de producción de datos que sigue requiriendo verificación externa.
\end{knowledgebox}

\subsection{Estrategia de reparación: depurar los datos como se depura software}

Las General Approaches incluyen limpiar o regularizar los datos, la recolección humana dirigida, generar anti-refusal synthetic data y el RL directo. Lo decisivo no es el nombre del método, sino localizar la causa raíz. A continuación, la ponente da la frase práctica más importante: \textbf{look at data like you'd debug software}. El function-calling refusal puede provenir de datos de self-knowledge; los rechazos ante documentos largos, de ejemplos del tipo «no tengo capacidad visual»; y la escritura creativa puede verse afectada por la proporción de datos de seguridad.

\lecturefigure{slide-42-refusal-general-approaches.jpg}{Rutas generales para reparar el over-refusal y resultados en WildChat/XSTest}{00:30:21--00:31:00}

\lecturefigure{slide-43-debug-data-like-software.jpg}{Revisar los datos como se depura software: rechazos distintos provienen de rutas distintas}{00:31:00--00:31:13}

\lecturefigure{slide-44-dataset-specific-refusals.jpg}{Posible acoplamiento entre los rechazos en escritura creativa y los datos de seguridad}{00:31:13--00:31:44}

\begin{importantbox}{Los cuatro objetos rastreables del data debugging}
\begin{enumerate}
  \item \textbf{Provenance}: si el ejemplo proviene de usuarios, anotadores, modelos o plantillas de políticas;
  \item \textbf{Slice}: a qué dominio de tarea, de riesgo y de idioma pertenece el ejemplo;
  \item \textbf{Influence}: qué slices de comportamiento cambian a la par después del entrenamiento;
  \item \textbf{Reversibility}: si es posible retirar esos datos o reducir su peso y volver a ejecutar experimentos locales.
\end{enumerate}
Una gran mezcla de datos sin provenance es difícil de depurar; «agregar algunos ejemplos buenos más» puede ocultar el conflicto en lugar de resolverlo.
\end{importantbox}

\teachervoice{00:28:35--00:31:44, la clase pasa de los product prompts, los synthetic prompts y XSTest a la root cause en los datos. La docente dice explícitamente que cada tipo de refusal puede deberse a un conjunto de datos distinto y que hay que revisarlo como cuando se depura software, en lugar de confiar en una perilla genérica de safety/helpfulness.}

\subsection{Helpfulness--Harmlessness es una optimización con restricciones, no un control deslizante único}

El data debugging anterior mostró que cada tipo de rechazo tiene una causa raíz distinta; ahora hay que resolver una decisión de producto que abarca todas las categorías: cuando el modelo está más dispuesto a responder, puede aumentar a la vez la información que infringe las políticas; cuando es más conservador, perjudica a los usuarios legítimos. Este problema no se resuelve con un control deslizante de «nivel de seguridad», porque el false refusal y el harmful compliance no son simétricos en costo, distribución ni frontera de responsabilidad. El objetivo puede escribirse como una optimización con restricciones:
\[
\max_{\pi}\;\mathbb{E}[H(\pi)]
\quad \text{s.t.}\quad
\mathbb{E}[U(\pi)]\le \epsilon,
\]
donde $H$ es la utilidad, $U$ es el riesgo inaceptable y $\epsilon$ es el presupuesto de riesgo que fija el producto. Un sistema real también debe establecer restricciones distintas por dominio y considerar el costo asimétrico de los false positive y los false negative.

\lecturefigure{slide-45-helpfulness-harmlessness.jpg}{Helpfulness y Harmlessness: buscar el equilibrio, no elegir un solo lado}{00:31:44--00:31:54}

\lecturefigure{slide-46-xstest-refusal-rates.jpg}{Incorrect refusal rate en XSTest: comparar versiones exige entender bien los ejemplos y la métrica}{00:31:54--00:32:10}

La figura de XSTest muestra diferencias notables en la incorrect refusal rate entre versiones de Claude. Al leer la figura, primero hay que confirmar que el eje vertical mide los rechazos erróneos y no la seguridad general, y después comparar las versiones del modelo; una barra baja respalda que hay menos rechazos ante benign prompts, pero no demuestra mayor seguridad ante harmful prompts ni representa todos los escenarios de idioma, herramientas o documentos largos. La cifra total de la figura debe interpretarse junto con las taxonomy slices.

\begin{warningbox}{Una mejora en el benchmark no significa que el problema del producto esté cerrado}
Los conjuntos públicos pueden no cubrir herramientas nuevas, contextos largos, system prompts reales ni la distribución de usuarios. La aceptación debería exigir: mejora en el benchmark público, que la regression interna no empeore, mejora en las slices de tráfico real, revisión humana de los casos graves y condiciones de monitoreo y reversión después del lanzamiento.
\end{warningbox}

\subsection{Vibe check: cómo los ejemplos cualitativos se convierten en evidencia formal}

Las tres pantallas before/after usan surveillance fiction, time travel government project y technology-free heist. Todas contienen palabras peligrosas en apariencia, pero son solicitudes creativas legítimas. Claude 2.1 tiende a rechazarlas; Claude 3 ofrece ayuda estructurada. Un ejemplo aislado no demuestra la capacidad general, pero estos ejemplos sirven para verificar si el objetivo de producto se cumplió «de manera significativa» y ayudan a las personas a detectar diferencias de tono que el benchmark no codifica.

\lecturefigure{slide-47-vibe-check-surveillance.jpg}{Vibe check I: la escritura de ciencia ficción sobre vigilancia pasa del rechazo a la ayuda estructurada}{00:32:10--00:32:13}

\lecturefigure{slide-48-vibe-check-time-travel.jpg}{Vibe check II: before/after de un proyecto secreto de viajes en el tiempo}{00:32:13--00:32:21}

\lecturefigure{slide-49-vibe-check-heist.jpg}{Vibe check III: before/after de un diálogo de heist sin tecnología moderna}{00:32:21--00:32:30}

\begin{knowledgebox}{Cuándo la qualitative eval no es «una ocurrencia sin fundamento»}
La evaluación cualitativa requiere fijar el prompt, la versión del modelo, la configuración de muestreo y la rubric; al menos dos reviewers juzgan de forma independiente y se registra el disagreement. Es especialmente útil para descubrir nuevos failure modes, problemas de tono y puntos de ruptura en la interacción. Una vez descubiertos, los ejemplos deben ampliarse a conjuntos de control, de variantes y contrafactuales antes de incorporarse a la regresión continua.
\end{knowledgebox}

\begin{lstlisting}[caption={Marco de ejecución por capas para evals de comportamiento}]
suites = [public_benchmarks, frozen_regressions, synthetic_slices,
          product_replays, qualitative_panels]
results = run_all(model, suites, fixed_system_and_sampling=True)
report = aggregate(results, by=[risk, task, language, tool, length])
gate(report.overall, report.worst_slice, severe_case_review)
archive(model_version, judge_version, prompts, traces, report)
\end{lstlisting}

\teachervoice{00:31:44--00:32:30, la ponente muestra primero el XSTest cuantitativo y después usa los tres pares before/after para hacer un vibe check. Nota de clase: los números se encargan de la regresión a escala y los ejemplos cualitativos, de juzgar si el comportamiento corresponde realmente a las convicciones del producto; ninguno sustituye al otro.}

\subsection{Resumen de la sección}

La behavior eval confiable parte del vector de comportamiento y, a través de la taxonomy, los mixed-source prompts, la provenance de los datos, las métricas con restricciones, los benchmarks y la qualitative review, conforma un protocolo de decisión. La lección más importante del caso de over-refusal no es «rechazar menos», sino que fallas distintas tienen causas raíz distintas. La siguiente sección lleva este enfoque al RL: cómo funciona el producto depende de cómo se construyen el entorno, las acciones, el reward y la verificación.
